\chapter{Feature Maps and Nonlinear Models}

\section{Basis Expansions}

Linear models are simple and computationally efficient, but their expressivity is limited: they can only represent linear relationships in the input space. A natural way to overcome this limitation is to transform the inputs before applying a linear model.

\medskip

\noindent
\textbf{Feature map.}  
A feature map is a function
\[
\phi : \mathcal{X} \to \mathbb{R}^m,
\]
which maps the original input into a new feature space.

\medskip

\noindent
We then define a model of the form
\[
f(x) = \langle w, \phi(x) \rangle.
\]

\medskip

\noindent
\textbf{Example (polynomial features).}  
For $x \in \mathbb{R}$,
\[
\phi(x) = (1, x, x^2, \dots, x^d).
\]

\medskip

\noindent
\textbf{Example (interaction features).}  
For $x \in \mathbb{R}^d$, include terms such as $x_i x_j$ to capture interactions.

\medskip

\noindent
\textbf{Key idea.}  
A linear model in feature space corresponds to a nonlinear model in the original input space.

\medskip

\noindent
\textbf{Interpretation.}  
Basis expansions allow us to approximate complex functions by combining simpler building blocks.

\medskip

\noindent
\textbf{Connection to approximation theory.}  
Polynomials, Fourier bases, and other function systems can approximate a wide range of functions when combined linearly.

\section{Linear Models in Feature Space}

Once a feature map is defined, learning proceeds as before:
\[
\min_w \; \frac{1}{n} \sum_{i=1}^n \ell(\langle w, \phi(x_i) \rangle, y_i).
\]

\medskip

\noindent
\textbf{Matrix formulation.}  
Let $\Phi \in \mathbb{R}^{n \times m}$ be the feature matrix:
\[
\Phi =
\begin{bmatrix}
- & \phi(x_1)^\top & - \\
- & \phi(x_2)^\top & - \\
& \vdots & \\
- & \phi(x_n)^\top & -
\end{bmatrix}.
\]

Then predictions are given by
\[
\hat{y} = \Phi w.
\]

\medskip

\noindent
\textbf{Insight.}  
All linear regression and classification methods extend directly to feature space by replacing $X$ with $\Phi$.

\medskip

\noindent
\textbf{Expressivity vs complexity.}
\begin{itemize}
    \item Increasing feature dimension $m$ increases expressivity
    \item But also increases risk of overfitting and computational cost
\end{itemize}

\medskip

\noindent
\textbf{Geometric viewpoint.}  
The data is embedded into a higher-dimensional space where it may become linearly separable.

\medskip

\noindent
\textbf{Example.}  
Points that are not separable by a line in $\mathbb{R}^2$ may become separable by a hyperplane after mapping into $\mathbb{R}^3$.

\section{The Kernel Idea (Intuition)}

A challenge arises when the feature space is very high-dimensional or even infinite-dimensional.

\medskip

\noindent
\textbf{Observation.}  
Many algorithms depend on inner products between feature vectors:
\[
\langle \phi(x), \phi(x') \rangle.
\]

\medskip

\noindent
\textbf{Kernel function.}  
A kernel is a function
\[
k(x, x') = \langle \phi(x), \phi(x') \rangle,
\]
which computes this inner product without explicitly evaluating $\phi(x)$.

\medskip

\noindent
\textbf{Key idea (kernel trick).}  
We can work implicitly in a high-dimensional feature space by using $k(x,x')$ directly.

\medskip

\noindent
\textbf{Examples of kernels.}
\begin{itemize}
    \item Polynomial kernel:
    \[
    k(x,x') = (\langle x, x' \rangle + c)^d
    \]

    \item Gaussian (RBF) kernel:
    \[
    k(x,x') = \exp\left(-\frac{\|x - x'\|^2}{2\sigma^2}\right)
    \]
\end{itemize}

\medskip

\noindent
\textbf{Interpretation.}
\begin{itemize}
    \item Kernels measure similarity between inputs
    \item They define an implicit feature space
\end{itemize}

\medskip

\noindent
\textbf{Practical advantage.}
\begin{itemize}
    \item Avoid explicit high-dimensional computations
    \item Enable flexible nonlinear models
\end{itemize}

\medskip

\noindent
\textbf{Limitation.}
\begin{itemize}
    \item Computational cost scales with number of data points
    \item Less scalable than modern deep learning approaches
\end{itemize}

\section{Expressivity through Features}

Feature maps determine what functions a model can represent.

\medskip

\noindent
\textbf{Key question.}  
How does the choice of $\phi$ affect the expressive power of the model?

\medskip

\noindent
\textbf{Insight.}  
A linear model in a sufficiently rich feature space can approximate highly complex functions.

\medskip

\noindent
\textbf{Tradeoff.}
\begin{itemize}
    \item Rich features $\rightarrow$ high expressivity
    \item But increased risk of overfitting
\end{itemize}

\medskip

\noindent
\textbf{Design principles for features.}
\begin{itemize}
    \item Encode domain knowledge (e.g., invariances)
    \item Capture relevant structure
    \item Maintain computational tractability
\end{itemize}

\medskip

\noindent
\textbf{From feature engineering to representation learning.}
\begin{itemize}
    \item Classical methods rely on manually designed features
    \item Deep learning learns features automatically from data
\end{itemize}

\medskip

\noindent
\textbf{Compositionality.}  
Complex features can be built by composing simpler transformations:
\[
\phi(x) = \phi_L(\phi_{L-1}(\cdots \phi_1(x))).
\]

\medskip

\noindent
This idea forms the foundation of deep neural networks.

\medskip

\noindent
\textbf{Geometric perspective.}  
Feature maps reshape the geometry of the data:
\begin{itemize}
    \item clusters become separable,
    \item nonlinear structures become linear,
    \item distances and similarities are redefined.
\end{itemize}

\section{A Unifying Perspective}

Feature maps provide a bridge between linear models and nonlinear behavior.

\medskip

\noindent
\textbf{Key components.}
\begin{itemize}
    \item Feature map $\phi$
    \item Linear model in feature space
    \item Optimization over parameters
\end{itemize}

\medskip

\noindent
\textbf{Two paradigms.}
\begin{itemize}
    \item \textbf{Fixed features:} designed manually (e.g., kernels)
    \item \textbf{Learned features:} optimized jointly with the model (deep learning)
\end{itemize}

\medskip

\noindent
\textbf{Core insight.}  
Nonlinearity in machine learning often arises not from the model itself, but from the representation of the data.

\section{Outlook}

This chapter introduced feature maps as a way to extend linear models:
\begin{itemize}
    \item basis expansions for nonlinear modeling,
    \item linear learning in transformed spaces,
    \item kernel methods for implicit high-dimensional features,
    \item the role of representation in expressivity.
\end{itemize}

\medskip

In the next chapter, we study support vector machines, which combine geometric insights with kernel methods to produce powerful classifiers.

\medskip

\noindent
\textbf{Guiding principle.}  
The choice of representation often determines the success of a learning system as much as the choice of model.